\chapter{Cronograma de Actividades}

\section{Linea de Tiempo}
(El impacto social hace referencia a los efectos que el proyecto investigado tiene sobre la población en general por lo tanto el estudiante deberá redactar el impacto social que espera tener su proyecto o investigación una vez  finalizada, es decir, que región geográfica o sector de la población se ven beneficiados.)

En esta sección listamos y explicamos las actividades planeadas en este proyecto. Además de explicar de forma detallada nuestra linea de tiempo para cada actividad. La Tabla \ref{fig:cronograma} muestra la relación de cada tarea con el tiempo en el que se ejecutará.

%\setlength{\tabcolsep}{6pt} % Default value: 6pt Horizoltal spacing
%\renewcommand{\arraystretch}{1}% Default value: 1 Vertical spacing

\section{Linea de Tiempo (segundo)}
El cronograma de actividades se muestra en la figura \ref{fig:cronograma}. 

\begin{figure}[!htb]
\caption{Diagrama de Gantt del cronograma de actividades.}
 \scalebox{0.8}{
  \begin{gantt}[xunitlength=1cm,fontsize=\small,titlefontsize=\small,drawledgerline=true]{9}{12}
    \begin{ganttitle}
      \titleelement{2016}{12}
    \end{ganttitle}
    \begin{ganttitle}
      \numtitle{1}{1}{12}{1}
    \end{ganttitle}
    \ganttbar[color=red]{\textbf{1. Revisión de la Literatura}}{1}{7}
    \ganttbar[color=red]{\textbf{2. Recolección de imágenes}}{4}{4}
    \ganttbar{\textbf{3. Presentación de plan de tesis}}{8}{1}
    \ganttbar{\textbf{4. Investigación de soluciones}}{8}{2}
    \ganttbar{\textbf{5. Implementación}}{10}{0.5}
    \ganttbar{\textbf{6. Evaluación}}{10.5}{0.5}
    \ganttbar{\textbf{7. Presentación de trabajo de tesis.}}{11}{1}    
  \end{gantt}
  }
  \label{fig:cronograma}
\end{figure}
